% #############################################################################
% This is Chapter 1
% !TEX root = ../main.tex
% #############################################################################
\fancychapter{Introduction}
\label{chap1:intro}

% #############################################################################

\section{Motivation}
\label{chap1:motivation}
Space is a harsh operating environment for electronics. Ionizing radiation can induce \acp{SEU} that flip bits in digital systems, corrupting architectural state and causing failures that are difficult to predict and reproduce. The problem is particularly relevant in \acp{FPGA}. Beyond \acp{FF} and memories, an \ac{FPGA} stores its logic and routing in configuration memory, which means that an \ac{SEU} may alter the circuit itself and corrupt behaviour until corrected. To protect architectures from such events, hardware \acp{FTM} are implemented at the \ac{RTL} level.

Developing and integrating \acp{FTM} introduces trade-offs in hardware cost, power, performance, and design complexity. Selecting the appropriate mechanism for a given architecture requires understanding these trade-offs through a comparative analysis across a representative set of \acp{FTM}. Some architectures implement multiple \acp{FTM}, which can serve as a partial comparison, but they are not designed as \ac{FTM}-testing platforms. Users are expected to use the provided mechanisms, with little support for adding new ones or evaluating them under the same conditions.

A rigorous comparison further depends on injecting faults in a controlled, repeatable way and quantifying the outcomes to measure fault coverage. The lack of reproducible, user-facing platforms that combine \ac{FTM} integration with controlled \ac{FI} and metric extraction is a major reason why broad comparative studies remain scarce.

The present thesis addresses that gap by developing a reproducible workflow to implement and evaluate \acp{FTM} on a configurable \ac{RISC-V} \ac{SoC} targeting \acp{FPGA}, enabling controlled \ac{FI} and quantitative fault coverage measurement alongside a comparative trade-off analysis of \acp{FTM}.


\section{Objective and Contributions}
\label{chap1:objectives}

The primary objective of this work was to develop \ac{Fatori-V} as a configurable \ac{RISC-V} platform for \ac{FTM} integration, controlled fault injection, and comparative evaluation. The following contributions were made.

\begin{enumerate}
    \item \textbf{Select and integrate the \ac{RISC-V} baseline.} The final baseline is IOb-SoC-Ibex \cite{iob-soc-ibex}, comprising the Ibex-to-\ac{AXI} bridge, Py2HWSW integration \cite{py2hwsw}, build-flow hooks for external Vivado Tcl and benchmark selection, native Ibex \acp{FTM}, and register wrapping.

    \item \textbf{Develop \ac{Fatori-V}.} The platform uses a single \texttt{.yaml} run file to configure hardware features, \acp{FTM}, benchmarks, \ac{FI} campaigns, generated files, and reported metrics. It automates validation, file generation, file movement, build execution, benchmark execution, logging, and result-table generation.

    \item \textbf{Implement \ac{FTM} evaluation hardware.} This includes Register Redundancy, Logic Redundancy, Self-Testing, the Fault Manager, the \ac{CSR}-based monitoring interface, configurable monitoring levels, and the error and correction counters used by the result flow.

    \item \textbf{Develop the unified \ac{FI} standalone tool.} The system supports \ac{SEM}-based \cite{SEM} configuration-memory injection, register and \ac{FF} injection through on-chip hardware, mixed target pools, reproducible target ordering, benchmark synchronization, reusable area and time profiles, and target-pool export and replay.

    \item \textbf{Add configuration memory targeting.} This includes \ac{ACME} \cite{ACME} essential-bit mapping for the Basys3 and XCKU040 boards, \ac{EBD}-based address generation, module and device injection scopes, and cache support for repeated \ac{ACME} expansions.

    \item \textbf{Develop the Pblock Assignment Model.} The model sizes hardware modules from equations empirically defined under different architectural configurations, and generates/assigns the pblocks used by the \ac{FI} system in configuration memory injections.

    \item \textbf{Integrate benchmarks and metrics.} This includes the common firmware wrapper, CoreMark \cite{coremark}, Embench-IoT \cite{embench-iot}, the Fatori Stress benchmark, pre and post metric sampling, benchmark-specific metric parsing, Vivado report parsing, and consolidated \texttt{.csv} and \texttt{.xlsx} result tables.

    \item \textbf{Validate \ac{Fatori-V}.} The platform validation covers Pblock Sizing Model accuracy and time-profile behaviour, using timestamped \ac{FI} campaigns.

    \item \textbf{Compare \acp{FTM}.} The final experiments compare native Ibex \acp{FTM}, Register M-of-N, Logic M-of-N, SecureIbex, the baseline, and the configuration with all \acp{FTM} enabled. The comparison uses common injection rates, register/configuration ratios, injection scopes, benchmarks, \ac{FT} outcomes, performance metrics, and hardware-cost metrics.
\end{enumerate}





% #############################################################################
\section{Document Organization}
The remainder of the document is organized as follows:
\begin{itemize}
    \item \textbf{Chapter 2} introduces the background concepts required for this project. It then presents the \ac{FT}/\ac{FI} testing platform state of the art and surveys \ac{SoC} candidate architectures for \ac{Fatori-V}'s baseline. An initial hardware backbone for \ac{Fatori-V} is chosen.

    \item \textbf{Chapter 3} presents the \ac{Fatori-V} platform design, its architecture layers, user interface, and experiment workflow.

    \item \textbf{Chapter 4} describes the initial attempt to build \ac{Fatori-V} on top of the VeeRwolf EL2 platform, including the planned \acp{FTM} and the lessons that motivated a change in direction.

    \item \textbf{Chapter 5} presents the IOb-SoC-Ibex baseline, the Ibex core and its native fault tolerance mechanisms, and the integration work required to support \ac{CPU} configuration and external control.

    \item \textbf{Chapter 6} details \ac{Fatori-V}'s Py2HWSW layer, which includes the support infrastructure, implemented \acp{FTM}, the fault manager, and the metrics monitoring system.

    \item \textbf{Chapter 7} details the \ac{Fatori-V} upper layer, explaining how experiments are defined, validated, transformed into build parameters, executed, and how results are delivered. The chapter also includes the benchmark porting, the Pblock Assigning Model, and \ac{FI} synchronization.

    \item \textbf{Chapter 8} details the standalone \ac{FI} tool, covering its architecture, the software controller, the on-chip injection hardware, and additional features.

    \item \textbf{Chapter 9} reports the platform validation experiments and the \ac{FTM} comparative study. It covers the experimental setup, validation results, computed metrics, obtained results, and a discussion on hardware cost, performance, and \ac{FT} behaviour.

    \item \textbf{Chapter 10} concludes the thesis by reflecting on how the work addresses the gap introduced earlier and proposing directions for future work.
\end{itemize}